Decentralized coordination of flexible loads through the alternating direction
method of multipliers (ADMM) is an established way to flatten the winter peaks
that cold-climate electric heating imposes on distribution feeders. Three gaps
limit its practical credibility. First, multi-agent schemes are fragile to
communication failures: when agents stop responding, the coordinated aggregate
degrades in ways rarely quantified. Second, the coordinated aggregate is almost
never checked against the physics of the feeder it is supposed to relieve. Third,
``comfort'' is usually a proxy on shifted energy rather than a check on modelled
indoor temperature. We close all three with a network-validated pipeline that
combines a comfort-aware sharing-ADMM coordinator, a machine-learning imputer
that reconstructs the heating of non-responsive agents, and an AC power-flow
validation stage on a synthetic North-American low-voltage residential feeder
behind an explicit distribution transformer. On a synthetic Trois-Rivi\`eres cold
day, the coordination flattens the aggregate peak while holding the worst modelled
indoor-temperature excursion to about $1\,^{\circ}\mathrm{C}$, roughly half that of
a naive energy-only flatten. Crucially, the flattened aggregate \emph{clears} a
transformer thermal overload and restores feeder feasibility rather than merely
lowering a statistical peak, and the cleared state stays feasible under heavy
communication failure before re-violating at a sharp knee. Comparing three
families of machine-learning regressors, we find that whether one imputes at all
matters far more than which method. The result is a fully reproducible study that
ties coordination quality, thermal comfort, and measurable feeder feasibility
together.
